\documentclass[11pt]{article}

\usepackage[margin=1in]{geometry}
\usepackage{amsmath,booktabs,array}
\usepackage[hidelinks]{hyperref}
\usepackage{xcolor}
\usepackage{enumitem}

\setlist[itemize]{leftmargin=1.5em,itemsep=0.25em,topsep=0.35em}
\setlist[enumerate]{leftmargin=1.8em,itemsep=0.25em,topsep=0.35em}
\newcommand{\code}[1]{\texttt{#1}}
\newcommand{\status}[1]{\textcolor{blue!65!black}{\textbf{#1}}}

\title{Racer-LLM Experiment Progress Log}
\author{}
\date{Last updated: August 12, 2026}

\begin{document}
\maketitle

\noindent
This is a living record of implementation decisions, experiments, and current
questions. It is intended for project updates rather than as finished paper
text. New work should be added as dated sections; earlier entries should remain
short and should only be revised when a factual error is found.

\section{Project direction}

The main goal is to study the behavior and limitations of LLM-based policies in
a real-time multi-agent racing task. The baseline work is meant to provide
stable and understandable comparisons, not to turn the project into an exact
reproduction of every low-level detail in earlier racing systems.

The current comparison will eventually include:
\begin{itemize}
  \item locally trained and released $\alpha$-RACER policies;
  \item an online iterated best-response (IBR) planner;
  \item a pure rule-based policy;
  \item a raw LLM policy;
  \item an LLM policy with explicit prompt rules.
\end{itemize}

The parameterized policies share the same local MPC and simulator wherever
possible. IBR is a separate online trajectory-optimization baseline and should
be described as such.

\section{Background notes used in the recent work}

\subsection{$\alpha$-RACER in brief}

$\alpha$-RACER represents each car's high-level policy using five parameters
\[
  \theta_i=(q_i,\zeta_i,s_{1,i},s_{2,i},s_{3,i}).
\]
The parameters control how the MPC follows and modifies a precomputed
single-agent raceline, including tracking aggressiveness, speed along the
raceline, overtaking, and blocking behavior.

For a joint policy $\theta=(\theta_0,\theta_1,\theta_2)$, car $i$ has a
long-term value
\[
  V_i(x,\theta)=\mathbb{E}\left[\sum_t\gamma^t r^i_t\right].
\]
The useful property of a potential function $\Phi$ is that a unilateral policy
change should produce a similar change in the potential and in the deviating
car's value:
\[
\begin{split}
  &\Phi(x,\theta'_i,\theta_{-i})-\Phi(x,\theta_i,\theta_{-i})\\
  &\qquad\approx
  V_i(x,\theta'_i,\theta_{-i})-V_i(x,\theta_i,\theta_{-i}).
\end{split}
\]
After this function is learned offline, the online policy searches for joint
parameters that give a high potential value and executes its own five-parameter
block.

\subsection{Wang-style IBR in brief}

The central idea of IBR is straightforward. Given the opponent's latest
predicted trajectory, the ego car optimizes its own trajectory. The opponent is
then optimized against the new ego trajectory. This alternation is repeated for
a small number of rounds:
\[
  \theta_i^{\ell+1}
  \in\arg\min_{\theta_i} f_i(\theta_i,\theta_{-i}^{\ell}).
\]

Wang's sensitivity-enhanced IBR (SE-IBR) additionally estimates how an
opponent's optimum will change in response to the ego trajectory. The response
is obtained by differentiating the opponent's KKT conditions. Intuitively, an
active collision constraint with a large multiplier indicates that changing the
ego trajectory may force the opponent to yield. This gives the ego more
information than ordinary IBR, which treats the current opponent trajectory as
fixed during each solve.

The original Wang implementation and our platform differ substantially in
trajectory parameterization, collision constraints, and dynamics. The recent
implementation therefore follows the alternating best-response idea using our
local model; it is not described as an exact reproduction.

\section{August 10--12, 2026: local $\alpha$-RACER training}

\status{Status: implemented, trained, and provisionally validated.}

\subsection{What was implemented}

The complete cached pipeline is now in
\code{train/alpha\_racer/alpha\_potential\_training.py}. It contains three
clearly separated stages.

\paragraph{1. Rollout collection.}
For each group, a baseline joint parameter set is sampled within the configured
bounds of $(q,\zeta,s_1,s_2,s_3)$. Four races are simulated:
\begin{enumerate}
  \item the baseline parameters for all three cars;
  \item a unilateral parameter change for car 0;
  \item a unilateral parameter change for car 1;
  \item a unilateral parameter change for car 2.
\end{enumerate}
Only one car's parameters change in each deviation race. This gives paired data
for learning the effect of unilateral policy changes.

At each control update, the code stores a 39-dimensional feature vector with
relative progress, vehicle states, track curvature, and all 15 policy
parameters. The local reward is currently
\[
  r^i_t=w_p\Delta s^i_t
  -w_o\mathbf{1}_{\mathrm{offroad},i}
  -w_c\mathbf{1}_{\mathrm{collision}}.
\]

\paragraph{2. Value-function training.}
After a rollout finishes, discounted returns are computed backwards:
\[
  G^i_t=\sum_{k=t}^{T-1}\gamma^{k-t}r^i_k.
\]
One neural value function is trained for each car using $G^i_t$ as its target.
This is Monte-Carlo learning: the full simulated rollout is collected before
the network is updated, rather than performing a temporal-difference update at
every simulator step.

\paragraph{3. Potential-function training.}
The three trained value networks produce targets for unilateral value changes.
The potential network is trained so that
\[
  \widehat\Phi(x,\theta')-\widehat\Phi(x,\theta)
  \approx
  \widehat V_i(x,\theta')-\widehat V_i(x,\theta).
\]
Both sides are evaluated at the same baseline state. This avoids mixing a policy
change with the different state trajectory that the policy later causes.

Existing rollout data and checkpoints are reused automatically. Rebuilding an
earlier stage invalidates and rebuilds the dependent later stages.

\subsection{Connection to the released code}

\begin{center}
\small
\begin{tabular}{p{0.20\linewidth}p{0.35\linewidth}p{0.35\linewidth}}
\toprule
Task & Released repository & Local project \\
\midrule
Collect races &
\code{collect\_data.bash}, \code{multi\_car\_blocking.py},
\code{merge\_data.py} &
\code{collect\_rollout\_data()} \\
Fit three values &
\code{train\_q\_model\_multi.py},
\code{train\_q\_model\_multi\_rel.py} &
\code{train\_value\_functions()} \\
Fit potential &
\code{train\_p\_model\_multi\_dist.py} &
\code{train\_potential\_function()} \\
Network definition &
\code{model\_arch.py::SimpleModel} &
\code{policy/alpha\_model.py::SimpleNeuralNetwork} \\
Online policy &
Released numerical simulator &
\code{policy/alpha\_racer.py::AlphaRacerPolicy} \\
\bottomrule
\end{tabular}
\end{center}

The local implementation keeps the original overall idea but replaces the ROS
simulation assumptions and implicit data ordering with the current project
simulator and explicit deviation groups.

\subsection{Training run recorded this week}

\begin{itemize}
  \item Track: Berlin.
  \item Vehicle model: local kinematic model.
  \item MPC horizon: 10.
  \item Data: 100 groups $\times$ 4 rollouts $\times$ 120 control samples.
  \item Rollout duration: 12 s.
  \item Discount: $\gamma=0.98$.
  \item Seed: 7.
  \item Value validation MSE: 21.43, 20.95, and 22.63.
  \item Best potential difference MSE: 0.0761.
\end{itemize}

The weights are stored under \code{train/weights/alpha\_racer/}. Released
weights are kept separately under \code{train/weights/alpha\_racer/upstream/}.
The runtime accepts both formats. All online searches use the canonical paper
parameters; released checkpoints receive their required input transformation
internally.

\subsection{Initial three-position checkpoint comparison}

As an initial check, the local checkpoint and the released myopic MPC checkpoint
were each run through the standard three-position sweep on Berlin. Every race
used three homogeneous $\alpha$-RACER cars with the same checkpoint and lasted
50 s. Thus each checkpoint produced three races and nine car trajectories.

\begin{center}
\small
\begin{tabular}{lcc}
\toprule
Checkpoint & Mean progress gain per car & Mean speed \\
\midrule
Released myopic MPC & 180.80 m & 3.56 m/s \\
Locally trained & 198.64 m & 3.90 m/s \\
\bottomrule
\end{tabular}
\end{center}

The local checkpoint improved mean progress gain by 17.85 m and mean speed by
0.35 m/s in this check. Neither configuration had a recorded collision. The
released checkpoint had no recorded off-road time; the local checkpoint had
0.08 s in total across the nine car trajectories.

Current interpretation: the locally trained checkpoint runs successfully and
is faster under the local dynamics, so it is suitable for use as a baseline.
This small three-position check is not a complete robustness study: the three
races mainly permute start order and do not cover independent random seeds,
different gaps, or different tracks.

\section{August 10--12, 2026: ordinary local IBR}

\status{Status: ordinary IBR is the default and is ready to freeze as a baseline.}

\subsection{What was implemented}

The main solver is \code{policy/ibr.py::\_WangIBRSolver}, selected by
\code{IBR\_IMPLEMENTATION=local\_wang}. Its state is
\[
  x_i=[s_i,e_i,\Delta\psi_i,v_i]^\top,
\]
and it directly optimizes throttle and steering over a finite horizon. The
objective rewards progress and penalizes centerline error, heading error, speed
error, control changes, and constraint slack. The main idea is
\[
\begin{split}
J_i \approx{}&-w_p(s_{i,N}-s_{i,0})\\
&+\sum_k\left[
w_e e_{i,k}^2+w_\psi\Delta\psi_{i,k}^2
+w_v(v_{i,k}-v_{\mathrm{rl}})^2
+w_{\Delta u}\|\Delta u_{i,k}\|^2\right]\\
&+\text{track slack penalty}+\text{collision slack penalty}.
\end{split}
\]

It uses the shared local vehicle parameters, friction-based steering/speed
constraint, track bounds, and a smooth superellipse collision constraint. The
simulator still uses oriented rectangles for physical contact detection.

The IBR planner follows the Frenet centerline geometrically. The precomputed
raceline supplies its target speed but not its lateral path. This is also the
broad behavior of the IBR released with $\alpha$-RACER. In contrast,
$\alpha$-RACER and the LLM-based parameter policy modify and track the geometric
raceline.

The earlier \code{historical\_strict} solver is retained as a historical local
port. It is not the primary baseline because it preserves brittle hard
constraints, an apparent collision-expression error from the released code, and
an ad hoc heading correction.

\subsection{Why SE is not used in the main baseline}

SE was implemented as an optional experiment, but the default sensitivity
weight is now zero.

The main issue is the interpretation of collision multipliers. Wang's SE term is
derived from the KKT conditions of the original hard-constrained trajectory
problem. Our local collision constraint has nonnegative slack with a quadratic
penalty. When collision slack is active, its multiplier is approximately
\[
  \lambda_c\approx2\rho_c\epsilon_c.
\]
The multiplier therefore depends strongly on the engineering choice of slack
penalty $\rho_c$. A sensitivity term using this multiplier can change scale when
we tune numerical feasibility, even if the strategic interaction is unchanged.

Tests with SE weights 0, 0.01, 0.05, and 0.10 produced almost identical final
progress. The calculated sensitivity loss was also very small. SE therefore
added theoretical and tuning complexity without a clear benefit in the current
platform. It remains optional exploratory code but is not planned for the main
comparison.

\subsection{Short validation run}

Three ordinary 10 s IBR races were run with each ego start position. Their final
progress values were permutations of
\[
  (46.175,39.850,33.830)\ \mathrm{m}.
\]
Across these runs:
\begin{itemize}
  \item no physical collision occurred;
  \item no car crossed the nominal $|e|>0.44$ m boundary;
  \item the $|e|>0.55$ m simulator recovery was never activated;
  \item all solver calls succeeded;
  \item maximum collision slack was about $1.4\times10^{-6}$.
\end{itemize}

A deliberately difficult test used 0.4 m start gaps and initial speeds
$(4.0,2.5,2.5)$ m/s. The initial clearance was almost too small for the 0.1 s
controller to avoid contact. The solver remained successful and all cars stayed
on track, but total physical collision time was 1.18 s and maximum collision
slack reached 0.91. This confirms that slack prevents solver collapse; it is not
a safety guarantee for an almost unavoidable initial condition.

Each 10 s simulation required approximately 32--36 s of wall-clock time on the
test machine. The current IBR can be used as an online optimization baseline,
but it should not yet be described as satisfying the real-time deadline.

\section{Current decisions}

\begin{itemize}
  \item Use the local $\alpha$-RACER checkpoint as the main learned
  non-LLM baseline, while retaining the released checkpoint for reference.
  \item Freeze \code{local\_wang} with sensitivity weight zero as the main IBR
  configuration.
  \item Keep SE outside the primary experiments unless a later result gives a
  clear reason to revisit the multiplier/slack coupling.
  \item Do not weaken IBR merely because the current LLM car is slower. First
  inspect the LLM-selected parameters, reference speed, and avoidance behavior.
  \item Record both the nominal track-boundary violation $|e|>0.44$ m and the
  simulator recovery activation $|e|>0.55$ m.
\end{itemize}

\section{Open questions and next experiments}

\begin{itemize}
  \item Why does the current LLM policy have lower average speed than IBR?
  Check its selected $(q,\zeta,s_1,s_2,s_3)$, raceline reference speed, lateral
  reference, and time spent reacting to opponents.
  \item Add or verify IBR track-slack diagnostics in saved summaries.
  \item Implement a pure rule-based baseline and restore explicit prompt rules
  as a separate LLM condition.
  \item Freeze a common scenario suite across tracks, start positions, gaps,
  speed perturbations, and lateral perturbations.
  \item Report computation latency and missed decision deadlines, not only race
  progress and safety outcomes.
  \item Decide whether the existing precomputed racelines are sufficient or
  whether all tracks should be regenerated with one fixed minimum-curvature and
  speed-profile pipeline.
\end{itemize}

\section{Future dated entries}

\subsection*{YYYY-MM-DD: short experiment name}

\status{Status: planned / running / completed / superseded.}

\begin{itemize}
  \item \textbf{Question:} What was this experiment intended to determine?
  \item \textbf{Configuration:} Policies, track, dynamics, initial conditions,
  checkpoint, seed, and important command-line options.
  \item \textbf{Result:} Main quantitative and qualitative observations.
  \item \textbf{Interpretation:} What can and cannot be concluded?
  \item \textbf{Decision:} What changes, if anything, for the next experiment?
\end{itemize}

\end{document}
